\documentclass{article}


% if you need to pass options to natbib, use, e.g.:
%     \PassOptionsToPackage{numbers, compress}{natbib}
% before loading neurips_2023


% ready for submission
%\usepackage{neurips_2023}


% to compile a preprint version, e.g., for submission to arXiv, add add the
% [preprint] option:
% \usepackage[preprint]{neurips_2023}


% to compile a camera-ready version, add the [final] option, e.g.:
     \usepackage[final]{neurips_2023}


% to avoid loading the natbib package, add option nonatbib:
%    \usepackage[nonatbib]{neurips_2023}


\usepackage[utf8]{inputenc} % allow utf-8 input
\usepackage[T1]{fontenc}    % use 8-bit T1 fonts
\usepackage{hyperref}       % hyperlinks
\usepackage{url}            % simple URL typesetting
\usepackage{booktabs}       % professional-quality tables
\usepackage{amsfonts}       % blackboard math symbols
\usepackage{pifont}
\usepackage{amssymb}
\usepackage{nicefrac}       % compact symbols for 1/2, etc.
\usepackage{microtype}      % microtypography
\usepackage{xcolor}         % colors
\usepackage{graphicx}
\usepackage{multirow}
\usepackage{comment}
\usepackage{mathtools}
\usepackage{graphicx}
\usepackage{caption}
\usepackage{subcaption}
\usepackage{tablefootnote}
\usepackage{textcomp} \usepackage{scalerel}
% \usepackage[outdir=./]{epstopdf}


\newcommand{\cmark}{\ding{51}}%
\newcommand{\xmark}{\ding{55}}%

\usepackage{caption}


\usepackage{hyperref}
%\hypersetup{
%  colorlinks   = true, %Colours links instead of ugly boxes
%  urlcolor     = magenta, %Colour for external hyperlinks
%  linkcolor    = magenta, %Colour of internal links
%  citecolor   = color5 %Colour of citations, could be ``red''
%}
\usepackage{url}
\usepackage{graphicx}
\usepackage{booktabs}
\usepackage{algorithm}
\usepackage{algpseudocode}
\usepackage{multirow}
\usepackage{enumitem}
\usepackage{cleveref}
\usepackage{xcolor,colortbl}
\usepackage{wrapfig,lipsum,booktabs}
\usepackage{bbm}
% \usepackage{subfigure}
\usepackage[most]{tcolorbox}
\usepackage{lmodern}
\usepackage{tikz}
\usetikzlibrary{arrows.meta} % For arrow tips

\hypersetup{colorlinks, linkcolor={red!55!black}, citecolor={blue!65!black}, urlcolor={blue!65!black}}

\renewcommand{\sectionautorefname}{Section}
\renewcommand{\subsectionautorefname}{Section}
\renewcommand{\subsubsectionautorefname}{Section}


\newenvironment{pcrfont}{\fontfamily{pcr}\selectfont}{\par}
\DeclareTextFontCommand{\textpcr}{\pcrfont}
\newenvironment{lmttfont}{\fontfamily{lmtt}\selectfont}{\par}
\title{Self-Rewarding Language Models}


\author{%
  Weizhe Yuan$^{1,2}$ \quad Richard Yuanzhe Pang$^{1,2}$  \quad {Kyunghyun Cho}$^{2}$  \\ {\bf Xian Li}$^{1}$ \quad {\bf Sainbayar Sukhbaatar}$^{1}$  \quad     \textbf{Jing Xu}$^{1}$ 
  \quad \textbf{Jason Weston}$^{1,2}$   \\\\
  $^1$ Meta ~~~~~~~~~~ $^{2}$ NYU \\
}


\begin{document}


\maketitle


\begin{abstract}
We posit that to achieve superhuman agents, future models require 
superhuman feedback in order to provide an adequate training signal.
Current approaches commonly train reward models from human preferences, which may then be
bottlenecked by human performance level,
and secondly these separate  frozen reward models  cannot then learn to improve during LLM training.
In this work, we study {\em Self-Rewarding Language Models}, where the language model itself is used via LLM-as-a-Judge prompting to provide its own rewards during training. We show that 
during Iterative DPO training that 
 not only does instruction following ability improve, but also the ability to provide high-quality rewards to itself. 
Fine-tuning Llama 2 70B on three iterations of our approach yields a model that outperforms 
many existing systems on the AlpacaEval 2.0 leaderboard, including Claude 2, Gemini Pro, and GPT-4 0613.
While there is much left still to explore,  this work opens the door to the possibility of models that can continually improve in both axes.
\end{abstract}




\section{Introduction}

Aligning Large Language Models (LLMs) 
using human preference data can vastly improve the instruction following 
performance of pretrained models \citep{ouyang2022training,bai2022training}. %,rafailov2023direct,xu2023some,}.
The standard approach of Reinforcement Learning from Human Feedback (RLHF) 
learns a reward model from these human preferences. The reward model is then frozen and used to train the LLM using RL, e.g., via  PPO \citep{schulman2017proximal}.  A recent alternative is to avoid training the reward model at all, and directly use human preferences to train the LLM, as in Direct Preference Optimization \citep[DPO;][]{rafailov2023direct}.
In both cases, the approach is bottlenecked by the size and quality of the human preference data, and in the case of RLHF the quality of the frozen reward model trained from them as well. 

In this work, we instead
%introduce
propose to train a self-improving reward model
%continually updating 
% that is continually updating during LLM alignment training, rather than being frozen, 
that, rather than being frozen,  is continually updating during LLM alignment, in order to avoid this bottleneck. 
The key to such an approach is to develop an agent that possesses all the abilities desired during training, rather than separating them out into distinct models such as a reward model and a language model. 
In the same way that pretraining and 
multitasking training of instruction following tasks
%general instruction following 
allow task transfer by training on many tasks at once 
% xcaruana1997multitask,
\citep{collobert2008unified,radford2019language,ouyang2022training}, incorporating the reward model into that same system allows task transfer between the reward modeling task and the instruction following tasks.

We thus introduce {\em Self-Rewarding Language Models}, that both (i) act as instruction following models generating responses for given prompts; and (ii) can generate and evaluate new instruction following examples to add to their own training set. We train these models using an Iterative DPO framework similar to that recently introduced in \citet{xu2023some}. Starting from a seed model,  in each iteration there is a process of {\em Self-Instruction creation} whereby candidate responses are generated by the model for newly created prompts, and  are then assigned rewards by that same model. The latter is implemented via LLM-as-a-Judge prompting, which can also be seen as an instruction following task. A preference dataset is built from the generated data, and the next iteration of the model is trained via DPO, see
 \autoref{fig:model}.

In our experiments,  we start with a Llama 2 70B \citep{touvron2023llama2} seed model fine-tuned on  Open Assistant \citep{kopf2023openassistant}, and then perform the above training scheme.
We find that not only does the instruction following performance improve from Self-Rewarding LLM alignment compared to the baseline seed model, but importantly
the reward modeling ability, which is no longer fixed, improves as well. 
This means that the model during iterative training is able, at a given iteration, to provide a higher quality preference dataset to itself than in the previous iteration.
While this effect likely saturates in real-world settings, it provides the intriguing possibility of obtaining reward models (and hence LLMs) that are superior to ones that could have been trained from the original human-authored seed data alone.



\section{Self-Rewarding Language Models}

Our approach first assumes access to a base pretrained language model, and a small amount of human-annotated seed data.
We then build a model that aims to possess two skills simultaneously:
%\begin{enumerate}[leftmargin=*]  % to save space
\begin{enumerate} %[topsep=0pt,itemsep=4pt,parsep=0pt]
\item {\em Instruction following}: given a prompt that describes a user request, the ability to generate a high quality, helpful (and harmless) response. 
\item {\em Self-Instruction creation}: the ability to generate and evaluate
%high quality
new instruction-following examples %(prompts and response pairs) 
to add to its own training set.
%new prompts that a user might plausibly ask, and to then also generate a high-quality response to them.
%\item {\em LLM-as-a-Judge}: The ability to judge the quality of instruction following, which we can use to judge the model's own generations, replacing the need for an external reward model.
\end{enumerate}

\if 0
%We then build a model that aims to be able to do three things simultaneously:
We then require three :
\begin{enumerate}
\item {\em Instruction following}: given a prompt that describes a user request, the ability to generate a high-quality, helpful (and harmless) response. 
\item {\em Self-Instruction creation}: the ability to generate new prompts that a user might plausibly ask, and to then also generate a high-quality response to them.
%\item {\em LLM-as-a-Judge}: The ability to judge the quality of instruction following, which we can use to judge the model's own generations, replacing the need for an external reward model.
\end{enumerate}
\fi

These skills are used so that the model can perform self-alignment, i.e., they are the components used to iteratively train itself using AI Feedback (AIF). % (sometimes referred to as RLAIF -- Reinforcement Learning from AI Feedback ).

{\em Self-instruction creation} consists of generating candidate responses and then the model itself judging their quality, i.e., it acts as its own reward model, replacing the need for an external one. This is implemented via the  {\em LLM-as-a-Judge} mechanism \citep{zheng2023judging}, i.e.,
by formulating the evaluation of responses as an instruction following task.
This self-created AIF preference data is used as a training set.



\begin{figure*}[t!]
    \centering
    \includegraphics[width=\linewidth]{figs/self_reward_main.pdf}
    \caption{ {\bf Self-Rewarding Language Models.} Our self-alignment method consists of two steps: (i) {\em Self-Instruction creation}:  newly created prompts are used to generate candidate responses from model $M_t$, which also predicts its own rewards via LLM-as-a-Judge prompting. (ii) {\em Instruction following training}: preference pairs are selected from the generated data, 
    which are used for training via DPO, resulting in model $M_{t+1}$. %This Iterative DPO procedure is repeated, 
    This whole procedure can then be iterated 
    resulting in both improved instruction following and reward modeling ability. }
    \label{fig:model}
\end{figure*}



Our overall {\em self-alignment} procedure is an iterative one, which proceeds by building a series of such models, with the aim that each improves over the last. 
Importantly, because the model can both improve its generation ability, and act as its own reward model through the same generation mechanism, this means the reward model itself can improve through these iterations, deviating from standard practices where the reward model is fixed \citep{ouyang2022training}.
We believe this can increase the ceiling of the potential for self-improvement of 
these learning models going forward, removing a constraining bottleneck.


We describe these steps in more detail below. An overview of the approach is illustrated in Figure \ref{fig:model}.
%An initial model is trained with seed instruction-following and LLM-as-a-judge data



\subsection{Initialization} \label{sec:init}

\paragraph{Seed instruction following data} We are given a seed set of human-authored 
(instruction prompt, response) general instruction following examples that we use for training in a supervised fine-tuning (SFT) manner, starting from a pretrained base language model. Subsequently this will be referred to as Instruction
Fine-Tuning (IFT) data.

\paragraph{Seed LLM-as-a-Judge instruction following data}
We also assume we are provided a seed set of (evaluation instruction prompt, evaluation result response) examples which can also be used for training. 
While this is not strictly necessary, as the model using IFT data will already be capable of training an LLM-as-a-Judge,  we show that such training data can give 
improved performance (see Appendix~\ref{app:self_reward_ift_only} for supporting results). %the results for starting from a model using IFT data only are shown in ).
In this data, the input prompt asks the model to evaluate the quality of a given response to a particular instruction. The provided evaluation result response consists of chain-of-thought reasoning (a justification), followed by a final score (in our experiments out of 5). The exact prompt format we chose is given  in \autoref{tab:eval_prompt}, which  instructs the LLM to evaluate the response using five additive criteria (relevance, coverage, usefulness, clarity and expertise), covering various aspects of quality. 
%This training data thus allows the LLM to perform the role of a reward model.
Subsequently this will be referred to as Evaluation Fine-Tuning (EFT) data.

%\footnote{ We also provide results of a Self-Rewarding  model trained using IFT data alone, which works worse than using IFT+EFT data, in Appendix~\ref{app:self_reward_ift_only}.}

%\subsection{Instruction following}

%Given human annotations of the quality of instruction following of given responses of the form (instruction, )
%A particular class of instructions are the set of prompts asking the model to evaluate the quality of a given response to a specific prompt 
%in a supervised fine-tuning (SFT) manner, starting from a pretrained base language model. 

We use  both these seed sets together during training. 

%We use both of these seed sets together to train the model for the first iteration (from $M_0$ to $M_1$; see \autoref{sec:model-sequence}). 

\subsection{Self-Instruction Creation} \label{sec:creation}

Using the model we have trained, we can make it self-modify its own training set.
%We can add this training data to the seed data.
%In our setup we explore adding high quality training data to the seed data.
%In our setup we explore adding generated training data to the seed data.
Specifically, we generate additional training data for the next iteration of training.


This consists of the following steps:


 \begin{enumerate}[topsep=0pt,itemsep=4pt,parsep=0pt] %[leftmargin=*]  % to save space

\item { Generate a new prompt}: We generate a new prompt $x_i$ using few-shot prompting, sampling prompts from the original seed IFT data, following the approach of \citet{wang2022self} and \citet{honovich2022unnatural}.\footnote{In our main experiments, responses and rewards, items (2) and (3), are generated by the model we have trained, but generating prompts is actually done by a model fixed in advance. However, we show that prompts can also be generated by the newly trained model in each iteration in Appendix~\ref{app:augmented_prompts_generation_using_newly_trained_models}.} % However, we believe the more general design choice, which we think should also work well, is the one described here. }

% while could generate prompts too, we used a given fixed pool, which we generated using chatllamam

\item {Generate candidate responses}: We then generate $N$ diverse candidate responses $\{y_i^1, \ldots, y_i^N\}$ for the given prompt $x_i$ from our model using sampling. %In our experiments we use $N=4$.

\item {Evaluate candidate responses}: Finally, we use the {LLM-as-a-Judge} ability of our same model to evaluate its own candidate responses with scores $r_i^n \in [0,5]$ (exact prompt given in \autoref{tab:eval_prompt}). 

 \end{enumerate}
 



\subsection{Instruction Following Training}

\if 0
\paragraph{Human seed data} We are given a seed set of human-annotated 
(instruction prompt, response) examples
that we use for training in a supervised fine-tuning (SFT) manner, starting from a pretrained base language model. 
\fi 

As previously described, training is initially performed with the seed IFT and EFT data (\autoref{sec:init}). 
This is then augmented with additional data via AI (Self-)Feedback. 
% In training iteration $\geq 2$, we use the data obtained via AI (self-)feedback. 
% using the Self-Instruction procedure of \autoref{sec:creation}. % during iterative training.

\paragraph{AI Feedback Training}
%In our self-alignment approach, a
After performing the self-instruction creation procedure, we can 
augment the  seed data with  additional examples for training, which we refer to as
AI Feedback Training (AIFT) data. 

To do this, we construct {\em preference pairs}, which are training data of the form (instruction prompt $x_i$, winning response $y_i^w$, losing response $y_i^l$).  To form the winning and losing pair we take the highest and lowest scoring responses from the $N$ evaluated candidate responses (see \autoref{sec:creation}), following \citet{xu2023some}, discarding the pair if their scores are the same. These pairs can be used for training with a preference tuning algorithm. We use DPO \citep{rafailov2023direct}. 
%in this inner part of the loop, where the entire training procedure is referred to as Iterative DPO.

\if 0
We try two variants of such feedback:
% \begin{itemize}[leftmargin=*]  % to save space

{\em Preference pairs}: We construct training data of the form (instruction prompt $x_i$, winning response $y_i^w$, losing response $y_i^l$).  To form the winning and losing pair we take the highest and lowest scoring responses from the $N$ evaluated candidate responses (see \autoref{sec:creation}), following \citet{xu2023some}, discarding the pair if their scores are the same. These pairs can be used for training with a preference tuning algorithm. We use DPO \citep{rafailov2023direct}.

{\em Positive examples only}: In this variant, we add additional examples of (instruction prompt, response) curated by the model to the seed set for supervised fine-tuning, following other approaches \citep{li2023self,adolphs2022cringe,gulcehre2023reinforced}, rather than constructing preference data. In this setup we only add examples where the candidate response was evaluated to give a perfect score of $r_i^n=5$. 
%(i.e., 5 out of 5). The training then uses the standard negative log-likelihood loss. 


While we report the results of both approaches in our experiments, we find that learning from  preference pairs gives superior performance, and hence recommend that approach.
\fi 

% \end{itemize}






\begin{figure}[t]
\centering
\small
\begin{tcolorbox}[colback=green!10!white, % Background color
                  colframe=green!30!white, % Frame color
                  width=0.99\textwidth, % Width of the tcolorbox
                  arc=4mm, % Radius of the rounded corners
                  auto outer arc,
                  ]
Review the user's question and the corresponding response using the additive 5-point scoring system described below. Points are accumulated based on the satisfaction of each criterion:\\
\\
- Add 1 point if the response is relevant and provides some information related to the user's inquiry, even if it is incomplete or contains some irrelevant content.\\
- Add another point if the response addresses a substantial portion of the user's question, but does not completely resolve the query or provide a direct answer.\\
- Award a third point if the response answers the basic elements of the user's question in a useful way, regardless of whether it seems to have been written by an AI Assistant or if it has elements typically found in blogs or search results.\\
- Grant a fourth point if the response is clearly written from an AI Assistant's perspective, addressing the user's question directly and comprehensively, and is well-organized and helpful, even if there is slight room for improvement in clarity, conciseness or focus.\\
- Bestow a fifth point for a response that is impeccably tailored to the user's question by an AI Assistant, without extraneous information, reflecting expert knowledge, and demonstrating a high-quality, engaging, and insightful answer.\\
\\
\\
User: \texttt{\color{red}<INSTRUCTION\_HERE>}\\
\\
<response>\texttt{\color{red}<RESPONSE\_HERE>}</response>\\
\\
After examining the user's instruction and the response:\\
\\
- Briefly justify your total score, up to 100 words.\\
- Conclude with the score using the format: ``Score: <total points>''\\
\\
Remember to assess from the AI Assistant perspective, utilizing web search knowledge as necessary. To evaluate the response in alignment with this additive scoring model, we'll systematically attribute points based on the outlined criteria.
% Please evaluate the following model response to the given instruction on a scale from 1-5.
%\\\\
%Original Instruction: [INSTRUCTION PROMPT]
%\\\\
%Model Response: [MODEL RESPONSE TO INSTRUCTION] 
\end{tcolorbox}
\caption{{\bf LLM-as-a-Judge prompt for our LLM to act as a  reward model} and provide self-rewards for its own model generations. The model is initially trained with seed training data of how to perform well at this task, and then improves at this task further through our self-rewarding  training procedure.}
\label{tab:eval_prompt}
\end{figure}




\subsection{Overall Self-Alignment Algorithm}

\paragraph{Iterative Training}
Our overall procedure trains a series of models $M_1,\dots,M_T$ where each successive model $t$ uses augmented training data created 
by the $t-1^\text{th}$ model.  We thus define AIFT($M_t$) to mean AI Feedback Training data created using model $M_t$. 

\paragraph{Model Sequence}
\label{sec:model-sequence}
We define the  models, and the training data they use as follows: 

\begin{itemize}[topsep=0pt,itemsep=4pt,parsep=0pt]
\item[$M_0$]: Base pretrained LLM with no fine-tuning.

\item[$M_1$]: Initialized with $M_0$, then fine-tuned on the IFT+EFT seed data using SFT.

\item[$M_2$]: Initialized with $M_1$, then trained with AIFT($M_1$) data using DPO.

\item[$M_3$]: Initialized with $M_2$, then trained with AIFT($M_2$) data using DPO.

\end{itemize}

This iterative training resembles the procedure used in Pairwise Cringe Optimization and specifically is termed  Iterative DPO,  introduced in \citet{xu2023some}; however, an external fixed reward model was used in that work.


\section{Experiments}

\subsection{Experimental Setup}

\paragraph{Base Model}
In our experiments we use Llama 2 70B \citep{touvron2023llama2} as our base pretrained model.

\subsubsection{Seed Training Data} \label{sec:seed_data}

\paragraph{IFT Seed Data}
We use the human-authored examples provided in the Open Assistant dataset \citep{kopf2023openassistant} for instruction fine-tuning. Following \citet{li2023self} we use 3200 examples,  by sampling  only first conversational turns in the English language that
are high-quality, based on their human annotated rank (choosing only the highest rank 0).
In our experiments, we compare to a model fine-tuned from the base model using only this data via supervised fine-tuning, and refer to it as our {\em SFT baseline}.

\paragraph{EFT Seed Data} The Open Assistant data also provides multiple ranked human responses per prompt  from which we can construct evaluation fine-tuning data. We split this into train and evaluation sets, and use it to create LLM-as-a-Judge data. This is done by placing it in the input format given in \autoref{tab:eval_prompt}, which consists of the scoring criteria description, and the given instruction and response to be evaluated.\footnote{Note, the prompt, derived from  \citet{li2023self}, mentions ``utilizing web search'', but our model is not actually capable of this action.} For training targets, chain-of-thought justifications and final scores out of 5 are not directly provided, so we use the SFT baseline to generate such output evaluations for each input, %for a given instruction, 
and accept them into the training set if the ranking of their scores agrees with the human rankings in the dataset. We resample the training set by discarding some of the data that receives the most common score so that the scores are not too skewed, as we observe many samples receive a score of 4. This results in 1,630 train and 541 evaluation examples (which do not overlap with the IFT data).

\subsubsection{Evaluation Metrics}

We evaluate the performance of our self-rewarding models in two axes:
 their ability to follow instructions, and their ability
 as a reward model (ability to evaluate responses).
  
\paragraph{Instruction Following}
We evaluate head-to-head performance between various models using GPT-4 \citep{achiam2023gpt} as an evaluator over 256 test prompts (which we refer to as IFT test data) derived from various sources following \citet{li2023self} using the AlpacaEval evaluation prompt \citep{alpaca_eval}. We try the prompt in both orders comparing pairwise, and if the GPT-4 evaluations disagree we count the result as a tie.
We also perform a similar evaluation with humans (authors).
We additionally report results in the  AlpacaEval 2.0 leaderboard format which is evaluated over 805  prompts, and compute the win rate against the baseline GPT-4 Turbo model based on GPT-4 judgments. Further, we report results on MT-Bench \citep{zheng2023judging} a set of challenging multi-turn questions in various categories from math and coding to roleplay and writing, which uses GPT-4 to grade the model responses out of 10.
Finally we also test the models on a set of 9 NLP benchmarks: ARC-Easy \citep{Clark2018ThinkYH}, ARC-Challenge \citep{Clark2018ThinkYH}, HellaSwag \citep{DBLP:conf/acl/ZellersHBFC19}, SIQA \citep{DBLP:journals/corr/abs-1904-09728}, PIQA \citep{Bisk2020}, GSM8K \citep{cobbe2021gsm8k}, MMLU \citep{DBLP:conf/iclr/HendrycksBBZMSS21}, OBQA \citep{OpenBookQA2018} 
and NQ \citep{47761}. 

\paragraph{Reward Modeling}
We evaluate the correlation with human rankings on the evaluation set we derived from the Open Assistant dataset,  as described in \autoref{sec:seed_data}.
Each instruction has on average 2.85 responses with given rankings. 
We can thus measure the {\em pairwise accuracy}, which is how many times the order of the ranking between any given pair agrees between the model's evaluation and the human ranking.
We also measure the {\em exact match} count, which is how often the total ordering is exactly the same for an instruction. We also report the Spearman correlation and Kendall's $\tau$. 
Finally, we report how often the responses that the model scores a perfect 5 out of 5 are rated as the highest ranked by humans.

\subsubsection{Training Details}

\paragraph{Instruction following training}
The training hyperparameters we use are as follows.
For SFT we use learning rate $5.5e{-6}$ which decays (cosine) to $1.1e{-6}$ at the end of training, batch size $16$ and dropout $0.1$. We only calculate the loss on target tokens instead of the full sequence.
For DPO we use learning rate $1e{-6}$ which decays to $1e{-7}$, batch size $16$, dropout $0.1$, and a $\beta$ value of 0.1.
We perform early stopping by saving a checkpoint every 200 steps and evaluating generations using Claude 2 \citep{claude2} on 253 validation examples derived from various sources following \citet{li2023self}. This is evaluated pairwise against the previous step's generations using the AlpacaEval evaluation prompt format \citep{alpaca_eval}.

\paragraph{Self-Instruction creation}
To generate new prompts we use a fixed model, Llama 2-Chat 70B with 8-shot prompting following Self-Instruct~\citep{wang2022self}, where we sample six demonstrations from the IFT data and two from the model generated data, and use decoding parameters T = 0.6, p = 0.9. We use their prompt template for non-classification tasks and apply the same filtering techniques, including the ROUGE-L~\citep{lin-2004-rouge} similarity check, keyword filtering, and length filtering. Except for the prompt generation part, the other parts of the creation pipeline (generating the response, and evaluating it) use the Self-Rewarding model being trained. For candidate response generation we sample $N=4$ candidate responses with temperature $T = 0.7$, $p=0.9$. 
When evaluating candidate responses, as there is variance to these scores, in our experiments we also use sampled decoding (with the same parameters) and generate these evaluations multiple (3) times and take the average. We added
3,964 such preference pairs to form the AIFT($M_1$) dataset used to train $M_2$ via DPO, 
and 6,942 pairs  to form AIFT($M_2$) used to train $M_3$.



\subsection{Results}
\label{subsec:results}
\subsubsection{Instruction Following Ability}

Head to head performance results are provided in \autoref{fig:h2h}.

\begin{figure}[t!]
    \centering
    %\includegraphics[width=1\textwidth,trim={0 0.75cm 0 0}]{figs/h2h.pdf}
    \includegraphics[width=0.8\textwidth,trim={0 -1cm 0  0}]{figs/h2h2.pdf}
    \includegraphics[width=0.8\textwidth,trim={0 0 0 0}]{figs/h2h3.pdf}
   \caption{
   \label{fig:h2h}
   {\bf Instruction following ability improves with Self-Training:} 
     We evaluate our models using head-to-head win rates on diverse prompts using GPT-4. % 
    The SFT Baseline is on par with Self-Rewarding Iteration 1 ($M_1$).
     However, Iteration 2 ($M_2$) outperforms both Iteration 1 ($M_1$) and the SFT Baseline. 
     Iteration 3 ($M_3$) gives further gains over Iteration 2 ($M_2$), outperforming 
     $M_1$, $M_2$ and the SFT Baseline by a large margin. 
 %%   %\caption{{\bf Head-to-head win rates of Self-Rewarding Models vs SFT Baseline}. Self-Rewarding Iteration 2 ($M_2$) outperforms Iteration 1 ($M_1$) and the SFT Baseline.  The SFT Baseline is on par with Iteration 1 ($M_1$). }
    }
%         Self-Rewarding Iteration 3 ($M_3$) gives improvements over Iteration 2 ($M_2$), Iteration 1 ($M_1$ and the SFT Baseline, outperforming each by a large margin. 
 %    While SFT Baseline is on par with Self-Rewarding Iteration 1 ($M_1$), 
\end{figure}
\paragraph{EFT+IFT seed training performs similarly to IFT alone}
We find that adding the Evaluation Fine-Tuning (EFT) task to training does not impact instruction following performance compared to using Instruction Fine-Tuning (IFT) data alone
%even though the data distributions are quite different (See Appendix.~\ref{app:ift_eft_distribution}), 
with an almost equal head to head (30.5\% wins vs. 30.9\% wins).
This is a positive result because it means the increased capability of a model to self-reward  does not affect its other skills. We can thus use IFT+EFT training as Iteration 1 ($M_1$) of our Self-Rewarding model, and then run further iterations.


\paragraph{Iteration 2 ($M_2$) improves over Iteration 1 ($M_1$) and SFT Baseline}
Iteration 2 of Self-Rewarding training  ($M_2$) provides superior instruction following to Iteration 1 ($M_1$)  with 55.5\% wins for $M_2$ compared to only 11.7\% for $M_1$ in a head to head evaluation. 
It provides similar gains  over the SFT Baseline as well (49.2\% wins vs. 14.5\% wins). Clearly, there is a large jump in performance from $M_1$ to $M_2$ by using the preference data AIFT($M_1$) provided by the reward model from Iteration 1.

\paragraph{Iteration 3 ($M_3$) improves over Iteration 2 ($M_2$)}
We see a further gain in Iteration 3 over Iteration 2, with 47.7\% wins for $M_3$ compared to only 12.5\% for $M_2$ in a head to head evaluation. Similarly, the win rate over the SFT Baseline for $M_3$ increases to 62.5\% wins vs. 9.8\%, i.e., winning more often than the $M_2$ model did. Overall, we see large gains from $M_2$ to $M_3$ through training using the preference data AIFT($M_2$) provided by the reward model from Iteration 2.

\paragraph{Self-Rewarding models perform well on AlpacaEval 2 leaderboard}
We evaluate our models on the AlpacaEval 2.0 leaderboard format, with results given in 
\autoref{tab:alpaca_leaderb}. We observe the same findings as in the head-to-head evaluations, that training iterations yield improved win rates, in this case over GPT4-Turbo, from 9.94\% in Iteration 1, to 15.38\% in Iteration 2, to 20.44\% in Iteration 3. Our Iteration 3 model outperforms many existing models in this metric, including Claude 2, Gemini Pro, and GPT4 0613. We show some selected models from the leaderboard in the table. We note that many of those competing models contain either proprietary alignment data (which is typically large, e.g., over 1M annotations in \citet{touvron2023llama2}) or use targets that are distilled from stronger models. In contrast, our Self-Rewarding model  starts from a small set of seed data from Open Assistant, and then generates targets and rewards from the model itself for further iterations of  training. 

\begin{table}[h]
    \caption{
    {\bf AlpacaEval 2.0 results} (win rate over GPT-4 Turbo evaluated by GPT-4). Self-Rewarding iterations yield improving  win rates. Iteration 3 ($M_3$) outperforms many existing models that use proprietary training data or targets distilled from stronger models.
    }
    \vspace{2mm}
  \label{tab:alpaca_leaderb}
  \centering

% \scalebox{1.0}{
  \begin{tabular}{llcc}
    \toprule
                   &                          & \multicolumn{2}{c}{ Alignment Targets} \\
    \textbf{Model} &    \textbf{Win Rate}  & Distilled & Proprietary\\
    %{\em Self-Rewarding (Llama 2 70B base)} \\ % + Open Assistant 3.2k)}\\
    \midrule  
    Self-Rewarding 70B \\
    ~~~~ {\em Iteration 1} ($M_1$) & 9.94\% & \\ % 1092
    ~~~~ {\em Iteration 2} ($M_2$) & 15.38\%&   \\ % 1552
    ~~~~ {\em Iteration 3} ($M_3$) & 20.44\% & \\ % 2552
    \midrule  
   %\multirow{3}{10em}{Self-Rewarding\\ (Llama 2 70B base)} & Self-Rewarding Llama2-70B-OA3.2k & \bf{83.71} \\
     \midrule 
   {\em Selected models from the leaderboard}\\
   GPT-4 0314	& 22.07\% &   & \cmark  \\ %	1371
   Mistral Medium & 	21.86\% &   & \cmark   \\ %1500
   Claude 2 &	17.19\%&   & \cmark  	\\ %1069
   Gemini Pro &	16.85\%  &   & \cmark  \\	%1315
   %Tulu 2+DPO 70B &	15.98\% \\	%	1418
   GPT-4 0613 &	15.76\% &   & \cmark  	\\	%1140
   %Claude 2.1 &	15.73\%	\\	%1096
   %Mistral 7B v0.2 &	14.72\% \\ %	1676
   GPT 3.5 Turbo 0613	&14.13\%  &   & \cmark  \\ %	1328
   LLaMA2 Chat 70B &	13.87\% &   & \cmark   \\ %	1790
   %Cohere Command &	12.90\%  \\ %	1983
   Vicuna 33B v1.3 &	12.71\% &   \cmark & \\ %	1479
   Humpback LLaMa2 70B &	10.12\% \\ %	1107
   %GPT 3.5 Turbo 0301 &	9.62\% &   & \cmark  \\ %	827
   %OpenChat V2-W 13B &	9.62\% \\ %	1566
   % \\ %Humpback LLaMa 65B &	9.43\% \\ %	1232
   %GPT 3.5 Turbo 1106	9.18\% \\ %	796
   % \\ %LLaMA2 Chat 13B &	7.70\% \\ %	1513
   %Vicuna 13B v1.3 &	7.14\%  &   \cmark & \\ %	1132
   Guanaco 65B &	6.86\% & & \\ %	1249
   Davinci001 &	2.76\%  &  & \cmark \\ %	296
   Alpaca 7B &	2.59\%  &   \cmark &   \\%	396
    \bottomrule
  \end{tabular}
 % }
\end{table}


\begin{figure}[h!]
    \centering
    \begin{subfigure}[b]{0.7\linewidth}
        \includegraphics[width=\columnwidth]{figs/topic_win_rate.pdf}
        % \caption{}
        \label{fig:topic_winrate}
    \end{subfigure}
    \captionsetup{skip=-10pt}  % Adjust this value to change the space
    \caption{AlpacaEval win rate breakdown for instruction categories (full names given in Appendix).
    Self-Rewarding models give gains across several topics, but tend to e.g. give less gains on mathematics and reasoning tasks.
    %We exclude clusters that contain fewer than 10 instructions.
    }
    \label{fig:fine_grained_winrate}
\end{figure}

% \paragraph{Preference optimization outperforms augmenting with positive examples only}
% As an ablation, we tried an alternative self-training procedure
% of adding high-quality self-instruction creation examples to supervised
% fine-tuning (without preference optimization), rather than DPO. Unfortunately, we could not find a configuration where this approach helped. For example, adding 11,254 such examples that scored 5 out of 5, and optimizing the mixing weight in training, still yielded a head to head with the SFT Baseline of 29\% wins vs 30\% wins, i.e., no improvement.




\paragraph{Fine-grained analysis}
As described earlier, the overall performance of the model in AlpacaEval improves with each iteration of training. It would be interesting to break down the overall performance improvement to see exactly what type of tasks these improvements come from. Therefore, we cluster the instructions in AlpacaEval test set into different groups based on three perspectives: (1) instruction category (2) instruction complexity (3) expected response length. We achieve this by using GPT-4. The detailed statistical information of the breakdown and the prompting techniques we used for getting this breakdown can be found in Appendix~\ref{app:alpacaeval_test_sample_clustering}. Results for  the instruction category are given in \autoref{fig:fine_grained_winrate}, and the other two in Appendix \autoref{fig:fine_grained_winrate_detailed_other_two}. From the results we can conclude that (i) Self-Rewarding models can substantially improve the win rate in most categories, but there are some tasks for which this approach does not improve, such as mathematics and logical reasoning, indicating that our current training approach mainly allows the models to better utilize their existing knowledge.
(ii) Through Self-Rewarding model training, the model's win rate increases on almost all tasks of different complexity, and especially on slightly more difficult tasks (complexity of 5, 6, 7 out of 10). (iii) The models also show a steady increase in the win rate on tasks with instructions with different expected response lengths.



\paragraph{Data distribution analysis}
We perform a t-SNE \citep{van2008visualizing} visualization of the IFT, EFT and AIFT($M_1$) data, shown in Appendix~\ref{app:ift_eft_distribution}. We find good overlap between the IFT and AIFT($M_1$) examples, which is desired, while the EFT examples lie in a different part of the embedding space, which can help explain why they would not affect IFT performance. 
We observe that generations from $M_1$ on AlpacaEval have an average length of 1092, for $M_2$ they are 
1552, and for $M_3$ they are 2552, so the model is learning to generate longer responses, which we note may be a factor in relative performance.





\begin{figure}[t!]
    \centering
    \includegraphics[width=0.8\textwidth,trim={0 0 0 0}]{figs/h2h_humaneval.pdf}
   \caption{
   \label{fig:h2h_humaneval}
   {\bf } 
     %We conduct human evaluation and evaluate our models using head-to-head win rates on sampled 50 instructions in IFT test. 
     {\bf Human evaluation results.} Iterations of Self-Rewarding ($M_1$, $M_2$ and $M_3$) provide progressively better head-to-head win rates compared to the SFT baseline, in agreement with  the automatic evaluation results. 
    }

\end{figure}

\begin{table}[h]
 %\small
% \setlength{\tabcolsep}{5.7pt}
    \caption{{\bf MT-Bench Results} (on a scale of 10). Self-Rewarding iterations yield improving scores across various categories. Math, code \& reasoning performance and iteration gains are smaller than for other categories, likely due to the makeup of the Open Assistant seed data we use.}
  \label{tab:mt-bench}
  \centering
\begin{tabular}{lccc}
\toprule
               &  Overall         & Math, Code &  Humanities, Extraction,  \\
               &  Score & \& Reasoning &  STEM, Roleplay \& Writing \\
                \midrule
{SFT Baseline}    & 6.85 & 3.93 & 8.60\\
{$M_1$}  & 6.78 & 3.83 & 8.55\\
{$M_2$}  & 7.01 & 4.05 & 8.79    \\                                                       
{$M_3$}   & 7.25 & 4.17 & 9.10  \\  
\bottomrule
\end{tabular}
\end{table}


\begin{table}[h!]
% \small
% \setlength{\tabcolsep}{6.0pt}
    \caption{
    {{\bf NLP Benchmarks.} Self-Rewarding models mostly tend to maintain performance compared to the  Llama 2 70B base model and the SFT Baseline, despite being fine-tuned on very different instruction-following prompts.}
    }
  \label{tab:core_knowledge_eval_results}
  \centering
  % \setlength{\tabcolsep}{4pt}
\begin{tabular}{lccccc}
\toprule
               & \multicolumn{1}{c}{{\begin{tabular}[c]{@{}c@{}}ARC $(\uparrow)$\\challenge\end{tabular}}}
               & \multicolumn{1}{c}{{\begin{tabular}[c]{@{}c@{}}HellaSwag\\$(\uparrow)$\end{tabular}}}
               & \multicolumn{1}{c}{{\begin{tabular}[c]{@{}c@{}}GSM8K\\$(\uparrow)$\end{tabular}}}
               %& \multicolumn{1}{c}{{GSM8K}} 
               %& \multicolumn{1}{c}{{\begin{tabular}[c]{@{}c@{}}gsm8k\\ (em)\end{tabular}}}
               & \multicolumn{1}{c}{{\begin{tabular}[c]{@{}c@{}}MMLU\\$(\uparrow)$\end{tabular}}}
               %& \multicolumn{1}{c}{{MMLU}} 
               %& \multicolumn{1}{c}{{\begin{tabular}[c]{@{}c@{}}mmlu\\ (macro\_avg/acc)\end{tabular}}}
              % & \multicolumn{1}{c}{{\begin{tabular}[c]{@{}c@{}}obqa\\ (acc\_comp)\end{tabular}}}
              & \multicolumn{1}{c}{{\begin{tabular}[c]{@{}c@{}}NQ\\$(\uparrow)$\end{tabular}}} \\
               %& \multicolumn{1}{c}{{NQ}}\\
              %& \multicolumn{1}{c}{{\begin{tabular}[c]{@{}c@{}}nq\\ (em)\end{tabular}}} \\
                \midrule
{Llama 2}                                  & 57.40                                        & 85.30                                                             & 56.80                                                                              & 68.90                                                                                                                                                                     & 25.30                                                                           \\
{SFT Baseline}                                   & 55.97                                       & 85.17                                                          & 50.72                                                                             & 69.76                                                                                                                                                           & 34.35                                                                          \\
{$M_1$}                                    & 57.51                                       & 84.99                                                           & 60.27                                                                             & 69.34                                                                                                                                                         & 35.48                                                                          \\
{$M_2$}                                  & 54.51                                       & 84.27                                                    & 59.29                                                                             & 69.31                                                                                                                                                               & 33.07                                                                          \\
{$M_3$}                                  & 53.13                                       & 83.29                                                      & 57.70                                                                              & 69.37                                                                                                                                                                & 31.86  \\
\bottomrule
\end{tabular}
\end{table}




\paragraph{Human evaluation}  
To examine whether human judgments align with automatic evaluation results, we conduct human evaluations that compare SFT baseline generations with the generations from each iteration of Self-Rewarding training, i.e., models $M_1$, $M_2$, and $M_3$. Specifically, we randomly select 50 instructions from the IFT test set. Each instruction corresponds to three pairs of generations (i.e., baseline vs. $M_1$, baseline vs. $M_2$, baseline vs. $M_3$). For each pair of generations, we assign them to three different annotators (blind evaluation performed by the authors) to make a pairwise judgment, and take a majority vote to decide which generation is better. The human evaluation results are shown in \autoref{fig:h2h_humaneval}. We find that Self-Rewarding models from later iterations show a larger advantage over the SFT baseline model, which is consistent with GPT-4's judgments, and demonstrates the effectiveness of our iterative training procedure.

\paragraph{MT-Bench performance further validates these results}  
We report performance on MT-Bench in \autoref{tab:mt-bench} for the SFT baseline and iterations of the Self-Rewarding model.  We again see improvements across the iterations of training from $M_1$ to $M_3$, from 6.78 (out of 10) up to 7.25, with larger relative gains in the humanities, STEM, roleplay, writing and extraction categories, and smaller gains in the math, code and reasoning categories.
We expect that the latter is due to the seed prompts we use from Open Assistant tending to underemphasize the reasoning-based tasks.
We note also that these improvements are in spite of our method using and constructing prompts that only involve a single turn, given the MT-Bench benchmark itself is a multi-turn evaluation.





\paragraph{Self-rewarding models did not lose ability on NLP Benchmarks} 
As shown in \autoref{tab:core_knowledge_eval_results}, the performance of most NLP benchmark tasks evaluated are roughly similar to the baselines, with further detailed results on more datasets given in Appendix \autoref{tab:core_knowledge_eval_results_detailed}  that follow the same pattern.
 We hypothesize that given that our training data (seed data and synthetically generated data) are based on the Open Assistant prompts which may not be especially relevant to skills needed in the \autoref{tab:core_knowledge_eval_results} tasks, it is expected that the task performance stays roughly similar, or may even drop. 
 For example, in InstructGPT training \citep{ouyang2022training}
 they found that ``during RLHF fine-tuning, we observe performance regressions compared
to GPT-3 on certain public NLP datasets'' which they refer to as an ``alignment tax.'' 
 A clear future direction is to extend the self-rewarding paradigm to these types of tasks, by relying not only on seed prompts from Open Assistant, but also on seed prompts found in a larger variety of datasets. 








\subsubsection{Reward Modeling Ability}

Reward modeling evaluation results are provided in 
\autoref{tab:reward_perf}.

\paragraph{EFT augmentation improves over SFT baseline}
Firstly, we find that adding Evaluation Fine-Tuning (EFT) data into training, which gives examples to the model of how to act as an LLM-as-a-Judge, naturally improves its performance compared to training with Instruction Fine-Tuning (IFT) data alone. IFT data covers a wide range of general instruction tasks, and so does endow the SFT Baseline with the ability to evaluate responses; however, EFT data gives more examples of this specific task.  We find improvements across all five metrics measured when using IFT+EFT vs. IFT alone, e.g., 
the pairwise accuracy agreement with humans increases from 65.1\% to 78.7\%. 
%,  and 5-best\% from 39.6\% to 41.5\%.%\footnote{We have tried various LLM-as-Judge prompts to decide the most effective one for our experiments. Some other less effective prompts are listed in the Appendix.~\ref{app:xian_eft_prompt}.}


\begin{table}[t]
\caption{{\bf Reward Modeling ability improves with Self-Training}:  We evaluate the LLM-as-a-Judge via various metrics which measure alignment with held-out human preference data. Self-Rewarding Iteration 2 (Model $M_2$), which is trained using the self-reward model derived from its previous iteration $M_1$ outperforms Iteration 1 ($M_1$), while $M_1$ itself outperforms a standard SFT baseline model trained on only Instruction Fine-Tuning (IFT) data. Iteration 3 (Model $M_3$) gives further improvements over Iteration 2.
}
\label{tab:reward_perf}
\vspace{2mm}
\begin{center}
% \begin{small}

% \begin{sc}
\begin{tabular}{lcccc}
\toprule
%Model & Pairwise accuracy $(\uparrow)$ &  Exact Match &  Spearman corr. &  ktau corr. & 5-best\% \\
                  &                    &  \multicolumn{3}{c}{Self-Rewarding Models}   \\ 
% \cline{3-5}
                  % &                    &  Iteration 1 & Iteration 2   & Iteration 3  \\ 
% \cline{3-5}
Model             &      SFT Baseline  &  Iter 1 ($M_1$) & Iter 2 ($M_2$)  & Iter 3 ($M_3$) \\
% $M_1$        & $M_2$   & $M_3$ \\
\midrule
{Training data}  &       {\footnotesize{IFT}}          &    {\footnotesize{IFT+EFT}}     &    {\footnotesize{IFT+EFT}}      &   {\footnotesize{IFT+EFT}{+AIFT($M_1$)}} \\
              &                    &               &   {\footnotesize {+AIFT($M_1$)}} &  {\footnotesize{+AIFT($M_2$)}}  \\
\midrule
Pairwise acc. $(\uparrow)$ & 65.1\% & 78.7\% & 80.4\% & 81.7\% \\
5-best \%    $(\uparrow)$ &   39.6\% & 41.5\% & 44.3\% & 43.2\% \\
Exact Match \%  $(\uparrow)$ &  10.1\% & 13.1\% & 14.3\% & 14.3\%  \\
Spearman corr. $(\uparrow)$ &  0.253 & 0.279 & 0.331 & 0.349 \\
Kendall $\tau$ corr.    $(\uparrow)$ &   0.233 & 0.253 & 0.315 & 0.324  \\
\bottomrule
\end{tabular}
% \end{sc}
% \end{small}
\end{center}
\end{table}


\paragraph{Reward Modeling ability improves with Self-Training}
We find that performing a round of self-reward training {\em improves the ability of the model at providing self-rewards for the next iteration}, in addition to its improved instruction following ability. Model $M_2$ (Iteration 2) is trained using the reward model from $M_1$ (Iteration 1), but provides improved performance on all five metrics compared to $M_1$.
For example,  pairwise accuracy  improves from 78.7\% to 80.4\%.
Iteration 3 ($M_3$) improves several of these metrics further compared to $M_2$, for example pairwise accuracy increases from 80.4\% to 81.7\%.
This performance gain is achieved despite there being no additional 
EFT data provided, and the examples created during the {\em Self-Instruction creation} loop do not tend to look like LLM-as-a-Judge training examples. We hypothesize that because the model is becoming better at general instruction following, it nevertheless also improves at
the  LLM-as-a-Judge task.

\paragraph{Importance of the LLM-as-a-Judge Prompt}
In these experiments we used the LLM-as-a-Judge prompt format shown  in  \autoref{tab:eval_prompt}.
In preliminary experiments we also tried various other prompts to decide the most effective one to use. For example, we tried the prompt proposed in \citet{li2023self} which also proposes a 5-point scale, but describes the options as multiple choice in a range of quality buckets, see Appendix \autoref{tab:eval_prompt_xian}.
 In contrast, our prompt describes the points as additive, covering various aspects of quality. We find a large difference between these two prompts when using the SFT Baseline, e.g. 65.1\% pairwise accuracy for ours, and only 26.6\% pairwise accuracy for theirs. See Appendix~\ref{app:xian_eft_prompt} for further details.
 


\section{Related Work}

Automatically improving or self-correcting large language models is becoming
a major focus of research. A recent survey from \citet{pan2023automatically}
attempts to summarize the topic. However, this is a rapidly moving area, and 
there are already promising new works not covered there.


\paragraph{Reinforcement Learning from Human Feedback (RLHF)}
Preference learning approaches such as in \citet{ziegler2019fine,stiennon2020learning,ouyang2022training,bai2022training}
train a fixed reward model from human preference data, and then use the reward model to train via reinforcement learning (RL), e.g. via Proximal Policy Optimization (PPO) \citep{schulman2017proximal}. Thus, the reward signal  in a certain sense  already comes from a model even in these works, but distilled from human data. Nevertheless, this is  commonly referred to as RL from Human Feedback (RLHF). 
Methods such as Direct Preference Optimization (DPO) \citep{rafailov2023direct} avoid training the reward model 
entirely, and instead directly train the LLM using human preferences. Several other such competing methods exist as well \citep{zhao2023slic,zheng2023click,yuan2023rrhf}, including Pairwise Cringe Optimization (PCO) \citep{xu2023some}. 
%, which was shown to outperform DPO and PPO on AlpacaFarm \citep{dubois2023alpacafarm}. 
PCO uses an iterative training approach similar to the one in our work, except with a fixed reward model, and that work also showed that Iterative DPO improves over DPO using the same scheme.
We note that other works have developed iterative preference training schemes
as well, e.g. \citet{adolphs2022cringe,gulcehre2023reinforced,xiong2023gibbs}.

\paragraph{Reinforcement Learning from AI Feedback (RLAIF)}
Constitutional AI \citep{bai2022constitutional}
uses an LLM to give feedback and refine responses, and uses this data to train a reward model. This fixed, separate reward model is then used  to train the language model via RL, called  ``RL from AI Feedback'' (RLAIF). 
\citet{lee2023rlaif} compare RLAIF and RLHF procedures and find the methods they compare perform roughly equally. They
use an  ``off-the-shelf'' LLM to perform LLM-as-a-Judge prompting to build a training set to train a fixed reward model, which is then used for RL training. They also experiment with using the fixed but separate LLM-as-a-Judge model directly, which the authors report is computationally expensive due to using it within PPO training (rather than the offline step in the iterative approach we use in our work, which is relatively computationally cheap).
Finally, SPIN \citep{chen2024self} recently showed they can avoid reward models entirely in an Iterative DPO-like framework by using human labels as the winning response in a pair, and the last iteration's generations as the losing response in the pair. The authors note this has the limitation that once the model generations reach human performance, they are bottlenecked. Further, each input prompt is required to have a human annotated response, in contrast to our work.


\paragraph{Improving LLMs via data augmentation (and curation)}
Several methods have improved LLMs by (self-)creating training data to augment fine-tuning. Self-Instruct \citep{wang2022self} is a method for self-instruction creation of prompts and responses, which can be used to improve a base LLM. We make use of a similar technique in our work, and then use our self-reward model to score them.
Several approaches have also created training data by
distilling from powerful LLMs, and shown 
a weaker LLM can then perform well. For example,  Alpaca \citep{alpaca} fine-tuned a Llama 7B  model with text-davinci-003 instructions created in the style of self-instruct.  Alpagasus \citep{chen2023alpagasus} employed a strong LLM-as-a-Judge (ChatGPT) to curate the Alpaca dataset and filter to a smaller set, obtaining improved results. Instruction Backtranslation \citep{li2023self} similarly augments and curates training data, but augmenting via backtranslating from web documents to predict prompts. The curation is done by the LLM(-as-a-Judge) itself, so can be seen as an instance of a self-rewarding model, but in a specialized setting.
Reinforced Self-Training (ReST) \citep{gulcehre2023reinforced} uses a fixed, external reward to curate new high-quality examples to iteratively add to the training set, improving performance. In our experiments, we found that adding only positive examples in a related manner did not help, whereas  preference pairs did help (see Appendix \autoref{app:positive_examples} for details).




\paragraph{LLM-as-a-Judge}
Using LLM-as-a-Judge prompting to evaluate language models has become a 
standard approach \citep{dubois2023alpacafarm,alpaca_eval,fernandes2023devil,bai2023benchmarking,saha2023branch}, and is being used to train reward models or curate data as well, as described above \citep{lee2023rlaif,chen2023alpagasus,li2023self}. While some works such as \citet{kim2023prometheus} create training data to train an LLM to perform well as a judge, to our knowledge it is not common to combine this training with general instruction following skills as in our work.



\section{Conclusion}

We have introduced Self-Rewarding Language Models, models capable of self-alignment via judging and  training on their own generations. The method learns in an iterative manner, where in each iteration the model creates its  own preference-based instruction training data. This is done by assigning rewards to its own generations via LLM-as-a-Judge prompting, and using Iterative DPO to train on the preferences. We showed that  this training both improves the instruction following capability of the model, as well as its reward-modeling ability across the iterations. While there are many avenues left unexplored, we believe this is exciting  because 
this means  the model is better able to assign rewards in future iterations for improving instruction following -- a kind of virtuous circle. 
While this improvement likely saturates in realistic scenarios,
it still allows for the possibility of continual improvement beyond the human preferences that are typically used to build reward models and instruction following models today.
%are typically trained with in current systems. % today. 






\section{Limitations}

While we have obtained  promising experimental results, we currently consider them preliminary because there are many avenues yet to explore, among them the topics of further evaluation, including safety evaluation, and understanding the limits of iterative training. 

We showed that the iterations of training improve both instruction following and reward modeling ability, but only ran three iterations in a single setting. A clear line of further research is to understand the ``scaling laws’’ of this effect both for more iterations, and with different language models with more or less capabilities in different settings.

% While we have evaluated our models using GPT-4 using head-to-head and AlpacaEval 2 leaderboard style evaluation, there are many other automatic evaluation benchmarks that one can  measure. Further, 
We observed an increase in length in model generations, and there is a known correlation between length and estimated quality, which is a topic that should be understood more deeply in general, and in our results in particular as well. It would also be good to understand if so-called ``reward-hacking’’ can happen within our framework, and in what circumstances. As we are using both a language model as the training reward, and a language model for final evaluation (GPT-4) in some of our benchmarks, even if they are different models, this may require a deeper analysis than we have provided. While the human evaluation we conducted did provide validation of the automatic results, further study could bring more insights.

Another clear further avenue of study is to conduct safety evaluations -- and to explore safety training within our framework.
Reward models have been built exclusively for safety in existing systems \citep{touvron2023llama2}, and a promising avenue here would be to use the LLM-as-a-Judge procedure to evaluate for safety specifically in our self-rewarding training process. 
Given that we have shown that reward modeling ability  improves over training iterations, this could mean that the safety of the model could potentially improve over time as well, with later iterations being able to catch and mitigate more challenging safety situations that earlier iterations cannot.


\newpage


\bibliography{sample}
\bibliographystyle{plainnat}




\newpage
\appendix
\onecolumn

\section{Appendix}
\subsection{Distributions of IFT, EFT and AIFT data}
\label{app:ift_eft_distribution}



\begin{figure}[h]
    \centering
    \begin{subfigure}[b]{0.49\textwidth}
        \includegraphics[width=\textwidth]{figs/instruction_distribution.pdf}
        \caption{Instruction distribution of IFT, EFT and AIFT data.}
        \label{fig:instruction_dist}
    \end{subfigure}
    % \hfill
    \begin{subfigure}[b]{0.49\textwidth}
        \includegraphics[width=\textwidth]{figs/response_distribution.pdf}
        \caption{Response distribution of IFT, EFT, and AIFT data.}
        \label{fig:response_dict}
    \end{subfigure}
    % \hfill
    % \begin{subfigure}[b]{0.32\textwidth}
    %     \includegraphics[width=\textwidth]{figs/full_text_distribution.pdf}
    %     \caption{Image 3}
    %     \label{fig:image3}
    % \end{subfigure}
    \caption{Distributions of both instructions and responses for IFT, EFT  and AIFT data.}
    \label{fig:ift_eft_distribution}
\end{figure}



We have plotted the distribution of instructions for IFT, EFT and AIFT($M_1$) data, and the distribution of responses for IFT, EFT and AIFT($M_1$) data in \autoref{fig:ift_eft_distribution}. It is clear that the IFT data and EFT data come from very different distributions while the IFT and AIFT($M_1$) data come from similar distributions.

\subsection{EFT Prompts}
\label{app:xian_eft_prompt}


The EFT prompt which we use in our main experiments is shown in \autoref{tab:eval_prompt}. 
\if 0
\begin{figure*}[t]
\centering
\small
\begin{tcolorbox}[colback=green!10!white, % Background color
                  colframe=green!30!white, % Frame color
                  width=0.99\linewidth, % Width of the tcolorbox
                  arc=4mm, % Radius of the rounded corners
                  auto outer arc,
                  ]
Review the user's question and the corresponding response using the additive 5-point scoring system described below. Points are accumulated based on the satisfaction of each criterion:\\
\\
- Add 1 point if the response is relevant and provides some information related to the user's inquiry, even if it is incomplete or contains some irrelevant content.\\
- Add another point if the response addresses a substantial portion of the user's question, but does not completely resolve the query or provide a direct answer.\\
- Award a third point if the response answers the basic elements of the user's question in a useful way, regardless of whether it seems to have been written by an AI Assistant or if it has elements typically found in blogs or search results.\\
- Grant a fourth point if the response is clearly written from an AI Assistant's perspective, addressing the user's question directly and comprehensively, and is well-organized and helpful, even if there is slight room for improvement in clarity, conciseness or focus.\\
- Bestow a fifth point for a response that is impeccably tailored to the user's question by an AI Assistant, without extraneous information, reflecting expert knowledge, and demonstrating a high-quality, engaging, and insightful answer.\\
\\
\\
User: \texttt{\color{red}<INSTRUCTION\_HERE>}\\
\\
\texttt{<}response\texttt{>}\texttt{\color{red}<RESPONSE\_HERE>}\texttt{<}/response\texttt{>}\\
\\
After examining the user's instruction and the response:\\
\\
- Briefly justify your total score, up to 100 words.\\
- Conclude with the score using the format: ``Score: \texttt{<}total points\texttt{>}''\\
\\
Remember to assess from the AI Assistant perspective, utilizing web search knowledge as necessary. To evaluate the response in alignment with this additive scoring model, we'll systematically attribute points based on the outlined criteria.
% Please evaluate the following model response to the given instruction on a scale from 1-5.
%\\\\
%Original Instruction: [INSTRUCTION PROMPT]
%\\\\
%Model Response: [MODEL RESPONSE TO INSTRUCTION] 
\end{tcolorbox}
\caption{{\bf LLM-as-a-Judge prompt for our LLM to act as a  reward model} and provide self-rewards for its own model generations. The model is initially trained with seed training data of how to perform well at this task, and then improves at this task further through our self-rewarding  training procedure.  (Note the prompt, derived from  \citet{li2023self}, states ``utilizing web search'', but our model is not actually capable of this.)}
\label{tab:eval_prompt}
\end{figure*}
\fi 

\paragraph{Other EFT prompts we have tried}

At first, we took the EFT prompt from \citet{li2023self} as shown in \autoref{tab:eval_prompt_xian}. However, we found that this prompt was not as effective as our additive score-counting prompt because the model needed to treat the task as a multiple-choice problem, and it was difficult for the model to break down this multiple-choice problem into sub-problems involving evaluating various aspects of the response. When using the model trained on 3,200 IFT data only, its performance on the EFT test set using our additive score-counting prompt and prompt from \citet{li2023self} is shown in \autoref{tab:compare_eft_prompt}.



\begin{figure}[t]
\centering
%\small
\begin{tcolorbox}[colback=green!10!white, % Background color
                  colframe=green!30!white, % Frame color
                  width=0.99\textwidth, % Width of the tcolorbox
                  arc=4mm, % Radius of the rounded corners
                  auto outer arc,
                  ]
Below is a question from an user and a candidate response. Please grade the response on a 5-point scale using the following criteria:\\
\\
1: It means the answer is incomplete, vague, off-topic, controversial, or not exactly what the user asked for. For example, some content seems missing, numbered list does not start from the beginning, the opening sentence repeats user's question. Or the response is from another person’s perspective with their personal experience (e.g. taken from blog posts), or looks like an answer from a forum. Or it contains promotional text, navigation text, or other irrelevant information. \\
2: It means the answer addresses most of the asks from the user. It does not directly address the user's question. For example, it only provides a high-level methodology instead of the exact solution to user's question. \\
3: It means the answer is helpful but not written by an AI Assistant. It addresses all the basic asks from the user. It is complete and self contained with the drawback that the response is not written from an AI assistant's perspective, but from other people's perspective. The content looks like an excerpt from a blog post, web page, or web search results. For example, it contains personal experience or opinion, mentions comments section, or share on social media, etc.\\
4: It means the answer is written from an AI assistant's perspective with a clear focus of addressing the instruction. It provide a complete, clear, and comprehensive response to user’s question or instruction without missing or irrelevant information. It is well organized, self-contained, and written in a helpful tone. It has minor room for improvement, e.g. more concise and focused.\\
5: It means it is a perfect answer from an AI Assistant. It has a clear focus on being a helpful AI Assistant, where the response looks like intentionally written to address the user's question or instruction without any irrelevant sentences. The answer provides high quality content, demonstrating expert knowledge in the area, is very well written, logical, easy-to-follow, engaging and insightful.\\
\\
\\
User: \texttt{\color{red}<INSTRUCTION\_HERE>}\\
\\
\texttt{<}response\texttt{>}\texttt{\color{red}<RESPONSE\_HERE>}\texttt{<}/response\texttt{>}\\
\\
Please first briefly describe your reasoning (in less than 100 words), and then write ``Score: \texttt{<}rating\texttt{>}'' in the last line. Answer in the style of an AI Assistant, with knowledge from web search if needed. To derive the final score based on the criteria, let's think step-by-step.
\end{tcolorbox}
\caption{LLM-as-a-Judge prompt taken from \citet{li2023self}.}
\label{tab:eval_prompt_xian}
\end{figure}

\begin{table}[h]
\begin{center}
\begin{small}
\begin{sc}
  \setlength\tabcolsep{20pt}
\begin{tabular}{lcc}

\toprule
%Model & Pairwise accuracy $(\uparrow)$ &  Exact Match &  Spearman corr. &  ktau corr. & 5-best\% \\

EFT Prompt             &   Multiple Choice prompt   &  Ours \\
\midrule
Pairwise accuracy $(\uparrow)$ & 26.6\% & 65.1\%\\
5-best \%    $(\uparrow)$ &  23.5\%  & 39.6\% \\
Exact Match \%  $(\uparrow)$ &  1.1\% & 10.1\% \\
Spearman corr. $(\uparrow)$ & -0.18 & 0.25 \\
Kendall $\tau$ corr.    $(\uparrow)$ & -0.16 & 0.23 \\
\bottomrule
\end{tabular}
\end{sc}
\end{small}
\end{center}
\caption{We tried various LLM-as-Judge prompts using the model trained with 3,200 IFT data only and found that our additive score-counting prompt worked best which demonstrates significant improvements in EFT performance comparing to the prompt used by \citet{li2023self}. 
\label{tab:compare_eft_prompt}
}
\end{table}
\subsection{Self-rewarding Models Using IFT Data Only}
\label{app:self_reward_ift_only}
To demonstrate the importance of the EFT data, we also trained a series of models starting with the model trained only on the IFT data. The following is the model sequence.

\begin{itemize}
\item[$M_0$]: Base pretrained LLM with no fine-tuning.
\item[$M_1^\prime$]: Initialized with $M_0$, then fine-tuned on the IFT seed data only using SFT.
\item[$M_2^\prime$]: Initialized with $M_1^\prime$, then trained with AIFT($M_1^\prime$) data using DPO.
\item[$M_3^\prime$]: Initialized with $M_2^\prime$, then trained with AIFT($M_2^\prime$) data using DPO.
\end{itemize}

Since we did not use EFT data to train the series of models, they were not always able to score the responses according to the format and even when they did, the scores given typically converged to 4. Therefore, even when starting from the same number of generated new prompts, we could only collect a very small number of valid training samples for DPO.
In total, we collected 541 pairs to form the AIFT($M_1^\prime$) dataset used to train $M_2^\prime$ via DPO, and 429 pairs to form AIFT($M_2^\prime$) used to train $M_3^\prime$. The win rates are shown in \autoref{fig:self_reward_from_sft_baseline}. From the figure we can conclude that EFT data helps to get better performance in the same number of iterations and the gap in performance between the model trained with EFT data and the model trained without EFT data widens in the later iterations.

\begin{figure*}[t]
    \centering
    \if 0
    \begin{subfigure}[b]{0.495\linewidth}
        \includegraphics[width=\columnwidth]{figs/h2h4.pdf}
        % \caption{}
        % \label{fig:topic_winrate}
    \end{subfigure}
    % \quad
    \begin{subfigure}[b]{0.495\linewidth}
        \includegraphics[width=\columnwidth]{figs/h2h5.pdf}
        % \caption{complexity.}
        % \caption{}
        % \label{fig:complexity_winrate}
    \end{subfigure}
    \fi
     \includegraphics[width=0.7\textwidth,trim={0 -1cm 0  0}]{figs/h2h4.pdf}
        \includegraphics[width=0.7\columnwidth]{figs/h2h5.pdf}
    \caption{\textbf{EFT data helps the self-rewarding loop:} We evaluated the series of models trained using self-reward loops starting from the model trained using only IFT data. We performed head-to-head win rates comparisons on the IFT test set. While $M_2^\prime$ can improve over the SFT baseline and $M_3^\prime$ can improve even more over the SFT baseline, they lag far behind the corresponding models ($M_2$, $M_3$) that started from a base model trained using both IFT and EFT data, see \autoref{fig:h2h}.}
    \label{fig:self_reward_from_sft_baseline}
\end{figure*}

%%%%%%%%%%%%%%%%%%%%%%%%%%%%%%%%%%%%%%%%%%%%%%%%%%%%%%%%%%%%





\begin{figure}[t]
\centering
\small
\begin{tcolorbox}[colback=green!10!white, % Background color
                  colframe=green!30!white, % Frame color
                  width=0.99\textwidth, % Width of the tcolorbox
                  arc=4mm, % Radius of the rounded corners
                  auto outer arc,
                  ]
\texttt{\color{red}<LIST ALL ALPACAEVAL INSTRUCTIONS>}\\                  
Given the above list of possible instructions, define  a maximum of 20 categories that would cover the types of intructions, for example recipes, reasoning tasks, general knowledge etc. Try to cover as many of the instructions as possible with the maximum 20 categories, while keeping the categories high-level, simple and easy to understand.
\end{tcolorbox}
\caption{Prompt used to obtain instruction categories on the AlpacaEval test set.}
\label{tab:instruction_categories}
\end{figure}


\begin{figure}[th!]
\centering
\small
\begin{tcolorbox}[colback=green!10!white, % Background color
                  colframe=green!30!white, % Frame color
                  width=0.99\textwidth, % Width of the tcolorbox
                  arc=4mm, % Radius of the rounded corners
                  auto outer arc,
                  ]
Instruction: \texttt{\color{red}<INSTRUCTION>}

Given the above, categorize it into one of the following 20 categories:\\
\texttt{\color{red}<LIST ALL CATEGORIES>}\\
\\
Secondly, score the instruction in terms of complexity: how complex you think it is to answer from 1-10 (where 10 is a complex question whereby first reasoning or breaking down the question into multiple subquestions for example might help improve the answer).\\
\\
Thirdly, indicate how long you think the response to the instruction should be, either (a) 1 sentence, (b) 1-3 sentences, (c) 1 paragraph, (d) 2 paragraphs, or (e) 3 or more paragraphs.\\
\\
\\
Provide your final response in the following format:\\
Category: \texttt{<}one of the 20 categories\texttt{>}\\
Complexity: \texttt{<}score out of 10\texttt{>}\\
Length: \texttt{<}length category\texttt{>}. Do not provide the actual response. 
\end{tcolorbox}
\caption{Prompt for categorizing instructions based on their topics, complexities and expected response lengths.}
\label{tab:instruction_category_complexity_expected_length_prompt}
\end{figure}





\begin{table}[t]
\centering
   \caption{
   \label{tab:instruction_category_breakdown} Breakdown of AlpacaEval test set instructions by instruction category.
   }
  % \small
\begin{tabular}{lrr}
\toprule
\textbf{Category}     & \textbf{Number} & \textbf{Percentage}  \\
\midrule
Science / Technology / Engineering                  & 134    & 16.65\%    \\
Professional / Business / Marketing                 & 77     & 9.57\%      \\
Social Interaction / Relationships / Human Behavior & 68     & 8.45\%      \\
Miscellaneous / Other                               & 61     & 7.58\%      \\
Mathematics / Logical Reasoning                     & 52     & 6.46\%      \\
Cooking / Recipes                                   & 48     & 5.96\%      \\
Software Development / Coding / Algorithms          & 44     & 5.47\%      \\
Travel / Geography / Exploration                    & 41     & 5.09\%      \\
Literature / Writing / Communication                & 39     & 4.84\%      \\
History / Social Studies                            & 38     & 4.72\%      \\
Entertainment / Media Analysis                      & 34     & 4.22\%      \\
Language Learning / Linguistics                     & 32     & 3.98\%      \\
Music / Audio / Arts                                & 30     & 3.73\%      \\
DIY Projects / Hobbies                              & 24     & 2.98\%      \\
Technology / Gadgets / Consumer Products            & 20     & 2.48\%      \\
Gaming / Game Development                           & 18     & 2.24\%      \\
Exercise / Health / Wellness                        & 16     & 1.99\%      \\
Philosophy / Ethics / Ideology                      & 15     & 1.86\%      \\
Sports / Athletics / Physical Activity              & 12     & 1.49\%      \\
Strategy / Problem-Solving / Critical Thinking      & 2      & 0.24\%      \\
\bottomrule
\end{tabular}
\end{table}


\begin{table}[t]
\centering
   \caption{
   \label{tab:instruction_complexity_breakdown} Breakdown of AlpacaEval test set instructions by instruction complexity. The instructions increase in complexity from 1 to 9, where 10 is a complex question that requires first reasoning or breaking the problem into sub-problems before it can be solved.
   }
%\small
\begin{tabular}{lrr}
\toprule
\textbf{Complexity}     & \textbf{Number} & \textbf{Percentage}  \\
\midrule
3 &  238 & 29.57\% \\
2 & 206 & 25.59\% \\
4 & 122 & 15.16\% \\
6 & 79 & 9.81\% \\
5 & 68 & 8.45\% \\
7 & 41 & 5.09\% \\
1 & 34 & 4.22\% \\
8 & 14 & 1.74\% \\
9 & 3 & 0.37\% \\
\bottomrule
\end{tabular}
\end{table}


\begin{table}[t]
\centering
   \caption{
\label{tab:instruction_expected_response_length_breakdown} Breakdown of AlpacaEval test set instructions by expected response length.
   }
\begin{tabular}{lrr}
\toprule
\textbf{Expected Length}     & \textbf{Number} & \textbf{Percentage}  \\
\midrule
1-3 sentences & 361 & 44.84\% \\
1 paragraph & 269 & 33.42\% \\
1 sentence & 143 & 17.76\% \\
2 paragraphs & 31 & 3.85\% \\
3 or more paragraphs & 1 & 0.13\% \\
\bottomrule
\end{tabular}
\end{table}


\begin{figure}[h!]
    \centering
    \hfill
    \begin{subfigure}[b]{0.47\linewidth}
        \includegraphics[width=\columnwidth]{figs/complexity_win_rate.pdf}
        % \caption{}
        \label{fig:complexity_winrate}
    \end{subfigure}
    \hfill
     \begin{subfigure}[b]{0.485\linewidth}
        \includegraphics[width=\columnwidth]{figs/expected_length_win_rate.pdf}
        % \caption{}
        \label{fig:expected_response_len_winrate}
    \end{subfigure}
    \hfill
    \captionsetup{skip=0pt}  % Adjust this value to change the space
    \caption{AlpacaEval win rate breakdown for instruction complexities (left) and expected response lengths (right).
    Self-Rewarding models give gains across most complexities and all response length ranges.
    %We exclude clusters that contain fewer than 10 instructions.
    }
    \label{fig:fine_grained_winrate_detailed_other_two}
\end{figure}





\subsection{Preference optimization outperforms augmenting with positive examples only}
\label{app:positive_examples}
We also tried an alternative self-training procedure
of adding high-quality self-instruction creation examples to supervised
fine-tuning (without preference optimization), rather than DPO. 
 In this variant, we add additional examples of (instruction prompt, response) curated by the model to the seed set for supervised fine-tuning, following other approaches \citep{li2023self,adolphs2022cringe,gulcehre2023reinforced}, rather than constructing preference data. In this setup we only add examples where the candidate response was evaluated to give a perfect score of $r_i^n=5$.  Unfortunately we could not find a configuration where this approach helped. For example, adding 11,254 such examples that scored 5 out of 5, and optimizing the mixing weight in training, still yielded a head to head with the SFT Baseline of 29\% wins vs 30\% wins, i.e., no improvement.

\subsection{Augmented Prompt Generation Using Newly Trained Models}
\label{app:augmented_prompts_generation_using_newly_trained_models}

In our experiments, for time efficiency, we have created a fixed pool of augmented prompts in advance using ChatLlama 70B. In a real interactive system, ideally, those prompts could come from real users so that we can ensure the models are trained to align with real user requirements. Here, we also examine whether  our newly trained Self-Rewarding models in each iteration can generate new prompts through in-context learning, instead of using ChatLlama 70B.
To check this, we constructed 30 prompts with in-context examples using the original seed IFT data as described in \autoref{sec:creation} and tested whether $M_1$, $M_2$ and $M_3$ still possess in-context learning ability and can generate high quality instructions. According to manual inspection, all models can generate novel instructions given in-context examples in all 30 cases. However, for $M2$ and $M3$, the model is likely to first generate a few instructions, then generate a separator, and then start responding to the instructions, so some postprocessing might be necessary.

\subsection{AlpacaEval Test Sample Clustering}
\label{app:alpacaeval_test_sample_clustering}
We used the GPT-4 (\texttt{gpt-4-1106-preview}) model to categorize the instructions in the AlpacaEval test set into clusters from three perspectives: (1) instruction category, (2) instruction complexity, and (3) expected response length. To obtain instruction categories for the AlpaceEval test set, we used the prompt in \autoref{tab:instruction_categories} and obtained 20 categories in total. Then, to cluster the instructions into different groups, we use the prompt in \autoref{tab:instruction_category_complexity_expected_length_prompt} for each test example. The corresponding statistics are given in \autoref{tab:instruction_category_breakdown}, \autoref{tab:instruction_complexity_breakdown}, \autoref{tab:instruction_expected_response_length_breakdown}. The fine-grained results on instruction complexity and expected response length are given in \autoref{fig:fine_grained_winrate_detailed_other_two}.

\newpage




\begin{table}[h]
\small
\setlength{\tabcolsep}{3.9pt}
    \caption{
    {{\bf NLP Benchmarks}. Self-Rewarding models mostly tend to maintain performance compared to the Llama 2 base model and the SFT Baseline, despite being fine-tuned on very different instruction-following prompts.}}
    \vspace{2mm}
  \label{tab:core_knowledge_eval_results_detailed}
  \centering
  \scalebox{0.86}{
\begin{tabular}{lccccccccc}
\toprule
                & \multicolumn{5}{c}{{\textbf{Commonsense Reasoning}}}                                                                                                                                                    & \multicolumn{1}{c}{{\textbf{Math}}}                                       & \multicolumn{3}{c}{{\textbf{World Knowledge}}}                                                                                                                                                                                                                            \\
                     \cmidrule(lr){2-6}   \cmidrule(lr){7-7} \cmidrule(lr){8-10} 
                & \multicolumn{1}{c}{{\begin{tabular}[c]{@{}c@{}}ARC\\easy\end{tabular}}} 
                & \multicolumn{1}{c}{{\begin{tabular}[c]{@{}c@{}}ARC\\challenge\end{tabular}}}
                & \multicolumn{1}{c}{{HellaSwag}} & \multicolumn{1}{c}{{SIQA}} & \multicolumn{1}{c}{{PIQA}} & \multicolumn{1}{c}{{\begin{tabular}[c]{@{}c@{}}GSM8K\\ (em)\end{tabular}}} & \multicolumn{1}{c}{{\begin{tabular}[c]{@{}c@{}}MMLU\\ (macro\_avg/acc)\end{tabular}}} & \multicolumn{1}{c}{{\begin{tabular}[c]{@{}c@{}}OBQA\\ (acc\_comp)\end{tabular}}} & \multicolumn{1}{c}{{\begin{tabular}[c]{@{}c@{}}NQ\\ (em)\end{tabular}}} \\
                \midrule
{Llama 2} & 80.20                                   & 57.40                                        & 85.30                                   & 50.70                              & 82.80                              & 56.80                                                                              & 68.90                                                                                         & 60.20                                                                                    & 25.30                                                                           \\
{SFT Baseline}     & 76.49                                  & 55.97                                       & 85.17                                  & 51.48                             & 82.59                             & 50.72                                                                             & 69.76                                                                                  & 57.80                                                                                    & 34.35                                                                          \\
{$M_1$}     & 78.14                                  & 57.51                                       & 84.99                                  & 53.02                             & 82.92                             & 60.27                                                                             & 69.34                                                                                 & 57.60                                                                                    & 35.48                                                                          \\
{$M_2$}     & 74.84                                  & 54.51                                       & 84.27                                  & 51.23                             & 81.94                             & 59.29                                                                             & 69.31                                                                                  & 57.60                                                                                    & 33.07                                                                          \\
{$M_3$}     & 72.35                                  & 53.13                                       & 83.29                                  & 49.28                             & 80.79                             & 57.70                                                                              & 69.37                                                                                  & 58.40                                                                                    & 31.86  \\
\bottomrule
\end{tabular}
}
\end{table}
\begin{table}[h]
\footnotesize
\setlength{\tabcolsep}{4pt}
    \caption{
    {\textbf{MT-Bench Fine-grained Results}. We list our models' performance on each problem category. Self-reward is especially effective in improving the model's ability in writing, role-playing, extraction, and STEM tasks.}}
    \vspace{2mm}
  \label{tab:mt_bench_detailed}
  \centering
   \scalebox{0.88}{
\begin{tabular}{lccccccccc}
\toprule
             & \textbf{Writing} & \textbf{Roleplay} & \textbf{Reasoning} & \textbf{Math} & \textbf{Coding} & \textbf{Extraction}        & \textbf{STEM}  & \textbf{Humanities} & \textbf{Overall}          \\
\midrule
SFT  & 8.83   & 8.15     & 5.30       & 3.00  & 3.50    & 6.90 & 9.18 & 9.95       & 6.85 \\
M1           & 9.10     & 7.65     & 4.35      & 3.05 & 4.10    & 7.20               & 8.93 & 9.85       & 6.78         \\
M2           & 9.10     & 8.00      & 4.60       & 3.30  & 4.25   & 7.65              & 9.40   & 9.80        & 7.01           \\
M3           & 9.58   & 8.73    & 4.80       & 3.50  & 4.20    & 7.80               & 9.45  & 9.95       & 7.25           \\

\bottomrule
\end{tabular}
}
\end{table}



\subsection{NLP Benchmark Results and MT-Bench Results}
We provide the detailed model performance on a number of NLP benchmarks in \autoref{tab:core_knowledge_eval_results_detailed} and on MT-Bench in \autoref{tab:mt_bench_detailed}. In particular, some NLP benchmarks including ARC-Challenge, HellaSwag, SIQA, PIQA, and OBQA are all text completion tasks. In these tasks, given the multiple choice options, we choose the option corresponding to the highest log probability scored by the models as the final answer. As such, the objective of these particular tasks is quite different from what our algorithm tries to optimize, so the results on these tasks may not reflect the true capability of our models.


\end{document}
\end{document}